\input{preamble/preamble.tex}

\geometry{margin=0.85in}
\AtBeginDocument{\small}

\newcommand{\ExamNameField}{\noindent\textbf{Name:}\ \rule{0.7\linewidth}{0.4pt}\par\medskip}

\begin{document}
	\ExamTitleBlock{10th grade}{Learning evidence 2.1: Analysis of Decisions (Practice Solutions)}{}
	
	\ExamSection{Problems}
	\begin{ExamProblems}
		\item
		\subsection*{C2}
		Decision alternatives are the membership models: Model A (standard), Model B (premium), and Model C (flex).
		The states of nature are the combined outcomes of demand (high, medium, low) and staffing costs (favorable, unfavorable).
		Each outcome represents the annual profit (in dollars) from membership fees minus staffing costs and fixed costs.
		Demand probabilities are $0.30$ (high), $0.45$ (medium), and $0.25$ (low), while staffing costs are $0.55$ (favorable) and $0.45$ (unfavorable), independent of demand.
		Therefore, the combined-state probabilities are products of the two probabilities.

		\subsection*{C3}
		Payoff = revenue $-$ cost, where revenue = (members $\times$ fee) and cost = (members $\times$ variable cost) $+$ fixed cost.
		\[
		\begin{aligned}
		\text{High, favorable: } &\Pi_A = 640(36-19)-6800=4080,\\
		&\Pi_B = 500(50-19)-9200=6300,\\
		&\Pi_C = 720(32-19)-5400=3960.\\[4pt]
		\text{High, unfavorable: } &\Pi_A = 640(36-27)-6800=-1040,\\
		&\Pi_B = 500(50-27)-9200=2300,\\
		&\Pi_C = 720(32-27)-5400=-1800.\\[4pt]
		\text{Medium, favorable: } &\Pi_A = 480(36-19)-6800=1360,\\
		&\Pi_B = 380(50-19)-9200=2580,\\
		&\Pi_C = 540(32-19)-5400=1620.\\[4pt]
		\text{Medium, unfavorable: } &\Pi_A = 480(36-27)-6800=-2480,\\
		&\Pi_B = 380(50-27)-9200=-460,\\
		&\Pi_C = 540(32-27)-5400=-2700.\\[4pt]
		\text{Low, favorable: } &\Pi_A = 320(36-19)-6800=-1360,\\
		&\Pi_B = 250(50-19)-9200=-1450,\\
		&\Pi_C = 380(32-19)-5400=-460.\\[4pt]
		\text{Low, unfavorable: } &\Pi_A = 320(36-27)-6800=-3920,\\
		&\Pi_B = 250(50-27)-9200=-3450,\\
		&\Pi_C = 380(32-27)-5400=-3500.
		\end{aligned}
		\]
		Combined probabilities: $0.30\cdot 0.55=0.165$, $0.30\cdot 0.45=0.135$, $0.45\cdot 0.55=0.2475$, $0.45\cdot 0.45=0.2025$, $0.25\cdot 0.55=0.1375$, $0.25\cdot 0.45=0.1125$.
		\begin{center}
			\begin{tabular}{l c c c c}
				\toprule
				State of nature & Probability & Model A & Model B & Model C \\
				\midrule
				High demand, favorable cost & 0.165 & 4080 & 6300 & 3960 \\
				High demand, unfavorable cost & 0.135 & -1040 & 2300 & -1800 \\
				Medium demand, favorable cost & 0.2475 & 1360 & 2580 & 1620 \\
				Medium demand, unfavorable cost & 0.2025 & -2480 & -460 & -2700 \\
				Low demand, favorable cost & 0.1375 & -1360 & -1450 & -460 \\
				Low demand, unfavorable cost & 0.1125 & -3920 & -3450 & -3500 \\
				\bottomrule
			\end{tabular}
		\end{center}

		\subsection*{C4}
		\[
		\begin{aligned}
		\max \text{Model A} &= 4080,\\
		\max \text{Model B} &= 6300,\\
		\max \text{Model C} &= 3960.
		\end{aligned}
		\]
		The Maximax choice is Model B with 6300. The criterion is \emph{risky} because it focuses only on the best possible payoff and ignores how likely it is.

		\subsection*{C5}
		\[
		\begin{aligned}
		\min \text{Model A} &= -3920,\\
		\min \text{Model B} &= -3450,\\
		\min \text{Model C} &= -3500.
		\end{aligned}
		\]
		The Maximin choice is Model B with $-3450$ because it has the least severe worst-case loss. This is \emph{safe} because it protects against the worst outcome.

		\subsection*{C1}
		Expected values use the combined-state probabilities.
		\[
		\begin{aligned}
		EV_A &= 0.165(4080)+0.135(-1040)+0.2475(1360)+0.2025(-2480)+0.1375(-1360)+0.1125(-3920)\\
		&= 673.2-140.4+336.6-502.2-187-441=-260.8,\\
		EV_B &= 0.165(6300)+0.135(2300)+0.2475(2580)+0.2025(-460)+0.1375(-1450)+0.1125(-3450)\\
		&= 1039.5+310.5+638.55-93.15-199.375-388.125=1307.9,\\
		EV_C &= 0.165(3960)+0.135(-1800)+0.2475(1620)+0.2025(-2700)+0.1375(-460)+0.1125(-3500)\\
		&= 653.4-243+400.95-546.75-63.25-393.75=-192.4.
		\end{aligned}
		\]
		The expected value criterion chooses Model B because $1307.9$ is the largest expected profit. This aligns with both the Maximax and Maximin outcomes for Model B, so it is the best overall decision under uncertainty.

		\item
		\subsection*{C2}
		Decision alternatives are the snack plans: Plan A (fresh-made items), Plan B (mixed menu), and Plan C (budget staples).
		The states of nature are demand levels (high, medium, low) with probabilities $0.30$, $0.50$, and $0.20$.
		Each outcome represents profit in hundreds of dollars after all costs.

		\subsection*{C3}
		Payoff = revenue $-$ cost, and the problem already provides the resulting profits for each plan and demand level.
		\[
		\begin{aligned}
		\text{High demand: } &\Pi_A=88,\ \Pi_B=74,\ \Pi_C=60,\\
		\text{Medium demand: } &\Pi_A=54,\ \Pi_B=66,\ \Pi_C=50,\\
		\text{Low demand: } &\Pi_A=-6,\ \Pi_B=28,\ \Pi_C=36.
		\end{aligned}
		\]
		\begin{center}
			\begin{tabular}{l c c c c}
				\toprule
				State of nature & Probability & Plan A & Plan B & Plan C \\
				\midrule
				High demand & 0.30 & 88 & 74 & 60 \\
				Medium demand & 0.50 & 54 & 66 & 50 \\
				Low demand & 0.20 & -6 & 28 & 36 \\
				\bottomrule
			\end{tabular}
		\end{center}

		\subsection*{C4}
		\[
		\begin{aligned}
		\max \text{Plan A} &= 88,\\
		\max \text{Plan B} &= 74,\\
		\max \text{Plan C} &= 60.
		\end{aligned}
		\]
		The Maximax choice is Plan A. It is \emph{risky} because it only considers the highest possible profit.

		\subsection*{C5}
		\[
		\begin{aligned}
		\min \text{Plan A} &= -6,\\
		\min \text{Plan B} &= 28,\\
		\min \text{Plan C} &= 36.
		\end{aligned}
		\]
		The Maximin choice is Plan C with 36 because it gives the best worst-case outcome. It is \emph{safe} because it prioritizes the least harmful result.

		\subsection*{C1}
		\[
		\begin{aligned}
		EV_A &= 0.30(88)+0.50(54)+0.20(-6)=26.4+27-1.2=52.2,\\
		EV_B &= 0.30(74)+0.50(66)+0.20(28)=22.2+33+5.6=60.8,\\
		EV_C &= 0.30(60)+0.50(50)+0.20(36)=18+25+7.2=50.2.
		\end{aligned}
		\]
		The expected value criterion chooses Plan B because $60.8$ is the largest expected profit. This differs from the Maximax choice (Plan A) and the Maximin choice (Plan C), so the expected value gives a balanced decision using the probabilities.
	\end{ExamProblems}
\end{document}
